% Options for packages loaded elsewhere
\PassOptionsToPackage{unicode}{hyperref}
\PassOptionsToPackage{hyphens}{url}
\documentclass[
]{article}
\usepackage{xcolor}
\usepackage[margin=22mm]{geometry}
\usepackage{amsmath,amssymb}
\setcounter{secnumdepth}{-\maxdimen} % remove section numbering
\usepackage{iftex}
\ifPDFTeX
  \usepackage[T1]{fontenc}
  \usepackage[utf8]{inputenc}
  \usepackage{textcomp} % provide euro and other symbols
\else % if luatex or xetex
  \usepackage{unicode-math} % this also loads fontspec
  \defaultfontfeatures{Scale=MatchLowercase}
  \defaultfontfeatures[\rmfamily]{Ligatures=TeX,Scale=1}
\fi
\usepackage{lmodern}
\ifPDFTeX\else
  % xetex/luatex font selection
\fi
% Use upquote if available, for straight quotes in verbatim environments
\IfFileExists{upquote.sty}{\usepackage{upquote}}{}
\IfFileExists{microtype.sty}{% use microtype if available
  \usepackage[]{microtype}
  \UseMicrotypeSet[protrusion]{basicmath} % disable protrusion for tt fonts
}{}
\makeatletter
\@ifundefined{KOMAClassName}{% if non-KOMA class
  \IfFileExists{parskip.sty}{%
    \usepackage{parskip}
  }{% else
    \setlength{\parindent}{0pt}
    \setlength{\parskip}{6pt plus 2pt minus 1pt}}
}{% if KOMA class
  \KOMAoptions{parskip=half}}
\makeatother
\usepackage{longtable,booktabs,array}
\usepackage{calc} % for calculating minipage widths
% Correct order of tables after \paragraph or \subparagraph
\usepackage{etoolbox}
\makeatletter
\patchcmd\longtable{\par}{\if@noskipsec\mbox{}\fi\par}{}{}
\makeatother
% Allow footnotes in longtable head/foot
\IfFileExists{footnotehyper.sty}{\usepackage{footnotehyper}}{\usepackage{footnote}}
\makesavenoteenv{longtable}
\setlength{\emergencystretch}{3em} % prevent overfull lines
\providecommand{\tightlist}{%
  \setlength{\itemsep}{0pt}\setlength{\parskip}{0pt}}
\usepackage{bookmark}
\IfFileExists{xurl.sty}{\usepackage{xurl}}{} % add URL line breaks if available
\urlstyle{same}
\hypersetup{
  hidelinks,
  pdfcreator={LaTeX via pandoc}}

\author{}
\date{}

\begin{document}

\section{Exact planar few-distance constructions with 111 and 81
points}\label{exact-planar-few-distance-constructions-with-111-and-81-points}

AI Agent Lab --- research draft, 4 October 2026

\textbf{Status.} Both existence statements below are exact and
independently checked. The independent adversarial literature audit
through 4 October 2026 is complete and found no prior equal or stronger
construction in the primary sources and author code it examined. These
are improvements over the named published lower bounds. This qualified
literature conclusion does not assert certainty about unpublished or
unindexed work, global optimality, or a new general geometric mechanism.

\subsection{Abstract}\label{abstract}

We give an explicit planar set of 111 points determining exactly 41
distances, together with an 81-point set determining exactly 31
distances. Consequently, for the unrestricted Euclidean-plane problem
with at most k distinct distances, G(31) \ensuremath{\ge} 81 and G(41) \ensuremath{\ge} 111. The
constructions exceed the corresponding published lower bounds 80 and 109
in Ahmed--Snevily (2013). The recent Bao--Yu (2025) table gives 79 at 31
and does not cover 41. Both witnesses are specified by short lists of
integer half-planes and complete row intervals in the triangular
lattice. A complete finite integer calculation proves their
cardinalities and distance counts without relying on an optimizer,
floating-point tolerance, or an optimality certificate. The 81-point
witness has one reflection symmetry; the 111-point witness has threefold
rotational and reflection symmetry. We document source comparisons,
independently checked artifacts, discovery procedures, and the limits of
the novelty claim.

\begin{figure}[ht]
\centering
\setlength{\unitlength}{1pt}
\begin{picture}(440,200)
{\color{red}
\qbezier(59.000,76.019)(107.000,107.196)(155.000,138.373)
\qbezier(59.000,76.019)(110.000,102.000)(161.000,127.981)
\qbezier(107.000,159.158)(110.000,102.000)(113.000,44.842)
\qbezier(119.000,159.158)(116.000,102.000)(113.000,44.842)
}
\put(59.000,117.588){\circle*{2.8}}
\put(65.000,127.981){\circle*{2.8}}
\put(71.000,138.373){\circle*{2.8}}
\put(59.000,96.804){\circle*{2.8}}
\put(65.000,107.196){\circle*{2.8}}
\put(71.000,117.588){\circle*{2.8}}
\put(77.000,127.981){\circle*{2.8}}
\put(83.000,138.373){\circle*{2.8}}
\put(89.000,148.765){\circle*{2.8}}
\put(59.000,76.019){\circle*{2.8}}
\put(65.000,86.412){\circle*{2.8}}
\put(71.000,96.804){\circle*{2.8}}
\put(77.000,107.196){\circle*{2.8}}
\put(83.000,117.588){\circle*{2.8}}
\put(89.000,127.981){\circle*{2.8}}
\put(95.000,138.373){\circle*{2.8}}
\put(101.000,148.765){\circle*{2.8}}
\put(107.000,159.158){\circle*{2.8}}
\put(71.000,76.019){\circle*{2.8}}
\put(77.000,86.412){\circle*{2.8}}
\put(83.000,96.804){\circle*{2.8}}
\put(89.000,107.196){\circle*{2.8}}
\put(95.000,117.588){\circle*{2.8}}
\put(101.000,127.981){\circle*{2.8}}
\put(107.000,138.373){\circle*{2.8}}
\put(113.000,148.765){\circle*{2.8}}
\put(119.000,159.158){\circle*{2.8}}
\put(77.000,65.627){\circle*{2.8}}
\put(83.000,76.019){\circle*{2.8}}
\put(89.000,86.412){\circle*{2.8}}
\put(95.000,96.804){\circle*{2.8}}
\put(101.000,107.196){\circle*{2.8}}
\put(107.000,117.588){\circle*{2.8}}
\put(113.000,127.981){\circle*{2.8}}
\put(119.000,138.373){\circle*{2.8}}
\put(125.000,148.765){\circle*{2.8}}
\put(89.000,65.627){\circle*{2.8}}
\put(95.000,76.019){\circle*{2.8}}
\put(101.000,86.412){\circle*{2.8}}
\put(107.000,96.804){\circle*{2.8}}
\put(113.000,107.196){\circle*{2.8}}
\put(119.000,117.588){\circle*{2.8}}
\put(125.000,127.981){\circle*{2.8}}
\put(131.000,138.373){\circle*{2.8}}
\put(137.000,148.765){\circle*{2.8}}
\put(95.000,55.235){\circle*{2.8}}
\put(101.000,65.627){\circle*{2.8}}
\put(107.000,76.019){\circle*{2.8}}
\put(113.000,86.412){\circle*{2.8}}
\put(119.000,96.804){\circle*{2.8}}
\put(125.000,107.196){\circle*{2.8}}
\put(131.000,117.588){\circle*{2.8}}
\put(137.000,127.981){\circle*{2.8}}
\put(143.000,138.373){\circle*{2.8}}
\put(107.000,55.235){\circle*{2.8}}
\put(113.000,65.627){\circle*{2.8}}
\put(119.000,76.019){\circle*{2.8}}
\put(125.000,86.412){\circle*{2.8}}
\put(131.000,96.804){\circle*{2.8}}
\put(137.000,107.196){\circle*{2.8}}
\put(143.000,117.588){\circle*{2.8}}
\put(149.000,127.981){\circle*{2.8}}
\put(155.000,138.373){\circle*{2.8}}
\put(113.000,44.842){\circle*{2.8}}
\put(119.000,55.235){\circle*{2.8}}
\put(125.000,65.627){\circle*{2.8}}
\put(131.000,76.019){\circle*{2.8}}
\put(137.000,86.412){\circle*{2.8}}
\put(143.000,96.804){\circle*{2.8}}
\put(149.000,107.196){\circle*{2.8}}
\put(155.000,117.588){\circle*{2.8}}
\put(161.000,127.981){\circle*{2.8}}
\put(131.000,55.235){\circle*{2.8}}
\put(137.000,65.627){\circle*{2.8}}
\put(143.000,76.019){\circle*{2.8}}
\put(149.000,86.412){\circle*{2.8}}
\put(155.000,96.804){\circle*{2.8}}
\put(161.000,107.196){\circle*{2.8}}
\put(149.000,65.627){\circle*{2.8}}
\put(155.000,76.019){\circle*{2.8}}
\put(161.000,86.412){\circle*{2.8}}
\put(110,185){\makebox(0,0){\small 81 points, 31 distances}}
\put(110,17){\makebox(0,0){\small Published comparison: 80}}
{\color{blue}
\qbezier(270.000,127.981)(330.000,96.804)(390.000,65.627)
\qbezier(276.000,138.373)(333.000,102.000)(390.000,65.627)
\qbezier(270.000,65.627)(327.000,102.000)(384.000,138.373)
\qbezier(270.000,65.627)(330.000,96.804)(390.000,127.981)
\qbezier(330.000,169.550)(327.000,102.000)(324.000,34.450)
\qbezier(330.000,169.550)(333.000,102.000)(336.000,34.450)
}
\put(270.000,127.981){\circle*{2.8}}
\put(276.000,138.373){\circle*{2.8}}
\put(270.000,107.196){\circle*{2.8}}
\put(276.000,117.588){\circle*{2.8}}
\put(282.000,127.981){\circle*{2.8}}
\put(288.000,138.373){\circle*{2.8}}
\put(294.000,148.765){\circle*{2.8}}
\put(270.000,86.412){\circle*{2.8}}
\put(276.000,96.804){\circle*{2.8}}
\put(282.000,107.196){\circle*{2.8}}
\put(288.000,117.588){\circle*{2.8}}
\put(294.000,127.981){\circle*{2.8}}
\put(300.000,138.373){\circle*{2.8}}
\put(306.000,148.765){\circle*{2.8}}
\put(312.000,159.158){\circle*{2.8}}
\put(270.000,65.627){\circle*{2.8}}
\put(276.000,76.019){\circle*{2.8}}
\put(282.000,86.412){\circle*{2.8}}
\put(288.000,96.804){\circle*{2.8}}
\put(294.000,107.196){\circle*{2.8}}
\put(300.000,117.588){\circle*{2.8}}
\put(306.000,127.981){\circle*{2.8}}
\put(312.000,138.373){\circle*{2.8}}
\put(318.000,148.765){\circle*{2.8}}
\put(324.000,159.158){\circle*{2.8}}
\put(330.000,169.550){\circle*{2.8}}
\put(282.000,65.627){\circle*{2.8}}
\put(288.000,76.019){\circle*{2.8}}
\put(294.000,86.412){\circle*{2.8}}
\put(300.000,96.804){\circle*{2.8}}
\put(306.000,107.196){\circle*{2.8}}
\put(312.000,117.588){\circle*{2.8}}
\put(318.000,127.981){\circle*{2.8}}
\put(324.000,138.373){\circle*{2.8}}
\put(330.000,148.765){\circle*{2.8}}
\put(336.000,159.158){\circle*{2.8}}
\put(288.000,55.235){\circle*{2.8}}
\put(294.000,65.627){\circle*{2.8}}
\put(300.000,76.019){\circle*{2.8}}
\put(306.000,86.412){\circle*{2.8}}
\put(312.000,96.804){\circle*{2.8}}
\put(318.000,107.196){\circle*{2.8}}
\put(324.000,117.588){\circle*{2.8}}
\put(330.000,127.981){\circle*{2.8}}
\put(336.000,138.373){\circle*{2.8}}
\put(342.000,148.765){\circle*{2.8}}
\put(348.000,159.158){\circle*{2.8}}
\put(300.000,55.235){\circle*{2.8}}
\put(306.000,65.627){\circle*{2.8}}
\put(312.000,76.019){\circle*{2.8}}
\put(318.000,86.412){\circle*{2.8}}
\put(324.000,96.804){\circle*{2.8}}
\put(330.000,107.196){\circle*{2.8}}
\put(336.000,117.588){\circle*{2.8}}
\put(342.000,127.981){\circle*{2.8}}
\put(348.000,138.373){\circle*{2.8}}
\put(354.000,148.765){\circle*{2.8}}
\put(306.000,44.842){\circle*{2.8}}
\put(312.000,55.235){\circle*{2.8}}
\put(318.000,65.627){\circle*{2.8}}
\put(324.000,76.019){\circle*{2.8}}
\put(330.000,86.412){\circle*{2.8}}
\put(336.000,96.804){\circle*{2.8}}
\put(342.000,107.196){\circle*{2.8}}
\put(348.000,117.588){\circle*{2.8}}
\put(354.000,127.981){\circle*{2.8}}
\put(360.000,138.373){\circle*{2.8}}
\put(366.000,148.765){\circle*{2.8}}
\put(318.000,44.842){\circle*{2.8}}
\put(324.000,55.235){\circle*{2.8}}
\put(330.000,65.627){\circle*{2.8}}
\put(336.000,76.019){\circle*{2.8}}
\put(342.000,86.412){\circle*{2.8}}
\put(348.000,96.804){\circle*{2.8}}
\put(354.000,107.196){\circle*{2.8}}
\put(360.000,117.588){\circle*{2.8}}
\put(366.000,127.981){\circle*{2.8}}
\put(372.000,138.373){\circle*{2.8}}
\put(324.000,34.450){\circle*{2.8}}
\put(330.000,44.842){\circle*{2.8}}
\put(336.000,55.235){\circle*{2.8}}
\put(342.000,65.627){\circle*{2.8}}
\put(348.000,76.019){\circle*{2.8}}
\put(354.000,86.412){\circle*{2.8}}
\put(360.000,96.804){\circle*{2.8}}
\put(366.000,107.196){\circle*{2.8}}
\put(372.000,117.588){\circle*{2.8}}
\put(378.000,127.981){\circle*{2.8}}
\put(384.000,138.373){\circle*{2.8}}
\put(336.000,34.450){\circle*{2.8}}
\put(342.000,44.842){\circle*{2.8}}
\put(348.000,55.235){\circle*{2.8}}
\put(354.000,65.627){\circle*{2.8}}
\put(360.000,76.019){\circle*{2.8}}
\put(366.000,86.412){\circle*{2.8}}
\put(372.000,96.804){\circle*{2.8}}
\put(378.000,107.196){\circle*{2.8}}
\put(384.000,117.588){\circle*{2.8}}
\put(390.000,127.981){\circle*{2.8}}
\put(354.000,44.842){\circle*{2.8}}
\put(360.000,55.235){\circle*{2.8}}
\put(366.000,65.627){\circle*{2.8}}
\put(372.000,76.019){\circle*{2.8}}
\put(378.000,86.412){\circle*{2.8}}
\put(384.000,96.804){\circle*{2.8}}
\put(390.000,107.196){\circle*{2.8}}
\put(372.000,55.235){\circle*{2.8}}
\put(378.000,65.627){\circle*{2.8}}
\put(384.000,76.019){\circle*{2.8}}
\put(390.000,86.412){\circle*{2.8}}
\put(390.000,65.627){\circle*{2.8}}
\put(330,185){\makebox(0,0){\small 111 points, 41 distances}}
\put(330,17){\makebox(0,0){\small Published comparison: 109}}
\end{picture}
\caption{Exact lattice constructions on the same drawing scale. Colored segments show all four pairs at squared distance 91 (left) and all six pairs at 127 (right). Comparison counts refer to Ahmed--Snevily (2013).}
\end{figure}

\subsection{1. Problem and statements}\label{problem-and-statements}

For a finite X \ensuremath{\subset} R\ensuremath{{}^2}, let D(X) = \{\ensuremath{\Vert}x-y\ensuremath{\Vert} : x,y \ensuremath{\in} X, x \ensuremath{\ne} y\}, and define

G(k) = max \{ \textbar X\textbar{} : X \ensuremath{\subset} R\ensuremath{{}^2} finite,
\textbar D(X)\textbar{} \ensuremath{\le} k \}.

The plane and Euclidean metric are unrestricted. There is no bounding
container, prescribed diameter, contact condition, symmetry restriction,
or lattice restriction in G. A lattice witness supplies a lower bound on
G; it supplies no upper bound and does not prove that optimal sets must
be lattice sets.

Write \ensuremath{\varphi}(a,b) = (a+b/2, \ensuremath{\surd}3 b/2) for a,b \ensuremath{\in} Z, and Q(u,v)=u\ensuremath{{}^2}+uv+v\ensuremath{{}^2}. Then

\ensuremath{\Vert}\ensuremath{\varphi}(a,b)-\ensuremath{\varphi}(c,d)\ensuremath{\Vert}\ensuremath{{}^2} = Q(a-c,b-d).

The map \ensuremath{\varphi} is injective. The quadratic form Q is positive for every
nonzero integer vector because 4Q(u,v)=(2u+v)\ensuremath{{}^2}+3v\ensuremath{{}^2}. Counting nonzero
squared distances is equivalent to counting distances.

\textbf{Theorem 1.} The set \ensuremath{\varphi}(S31) defined in Section 3 has 81 points
and exactly 31 distances. Therefore G(31) \ensuremath{\ge} 81.

\textbf{Theorem 2.} The set \ensuremath{\varphi}(S41) defined in Section 3 has 111 points
and exactly 41 distances. Therefore G(41) \ensuremath{\ge} 111.

These are existence statements with full certificates. Neither theorem
asserts G(31)=81 or G(41)=111. Monotonicity also gives G(k)\ensuremath{\ge}81 for k\ensuremath{\ge}31
and G(k)\ensuremath{\ge}111 for k\ensuremath{\ge}41, but these consequences need not improve
previously known bounds. In particular, Ahmed--Snevily already gives 115
at k=42: this manuscript claims \textbf{no improvement at 42}.

\subsection{2. Prior work and comparison
scope}\label{prior-work-and-comparison-scope}

Erdős and Fishburn formulated and studied maximum planar sets
determining k distances {[}1{]}. Ahmed and Snevily studied symmetric
lattice polygons and arbitrary orientations of attached triangular
blades; their main construction figure on printed page 5 explicitly
states g(31)\ensuremath{\ge}80, g(41)\ensuremath{\ge}109, and g(42)\ensuremath{\ge}115 {[}2{]}. Their g uses exactly
k distances. An exactly-k construction is valid for our at-most-k
definition; historical bounds must also be propagated monotonically
before comparison. The audited table through 50 retains 80 and 109 at
our two target indices.

Ahmed and Wildstrom proved the boundary-difference realization property
for cycle-bounded triangular-lattice sets {[}3{]}. We therefore make no
claim to have discovered boundary reduction. Balaji et al.~revisited
lattice constructions, corrected distance tables, and compared
structured arrays and lattice disks {[}4{]}. The initial preprint
appeared in 2019. The six-author journal publication appeared in 2020;
the seven-author 2023 arXiv revision also includes Lo Phillips and is
identified separately in the references. Szöllősi and Östergård's 2020
graph/Gröbner classification work concerns, among other cases, six
distances in the plane {[}5{]}; it is not a classification of our
large-k cases.

Bao and Yu use first-distance palettes and equiangular hexagons {[}6{]}.
Their displayed table at m=31 gives 79; comparison solely to that recent
table would understate the older 80-point construction. Their
first-31-distance restricted computation does not imply an upper bound
on arbitrary triangular-lattice palettes. Allowing larger distances,
observing multiplicities, and deleting points to remove distance classes
are all prior ideas: the numerical witnesses here are not evidence of a
newly discovered palette-exchange mechanism.

\begin{longtable}[]{@{}
  >{\raggedleft\arraybackslash}p{(\linewidth - 8\tabcolsep) * \real{0.2000}}
  >{\raggedleft\arraybackslash}p{(\linewidth - 8\tabcolsep) * \real{0.2000}}
  >{\raggedleft\arraybackslash}p{(\linewidth - 8\tabcolsep) * \real{0.2000}}
  >{\raggedleft\arraybackslash}p{(\linewidth - 8\tabcolsep) * \real{0.2000}}
  >{\raggedleft\arraybackslash}p{(\linewidth - 8\tabcolsep) * \real{0.2000}}@{}}
\toprule\noalign{}
\begin{minipage}[b]{\linewidth}\raggedleft
k
\end{minipage} & \begin{minipage}[b]{\linewidth}\raggedleft
Ahmed--Snevily 2013, exactly-k witness
\end{minipage} & \begin{minipage}[b]{\linewidth}\raggedleft
Bao--Yu 2025 displayed lower bound
\end{minipage} & \begin{minipage}[b]{\linewidth}\raggedleft
this manuscript, exactly-k witness
\end{minipage} & \begin{minipage}[b]{\linewidth}\raggedleft
improvement over audited comparison
\end{minipage} \\
\midrule\noalign{}
\endhead
\bottomrule\noalign{}
\endlastfoot
31 & 80 & 79 & 81 & +1 \\
41 & 109 & not tabulated & 111 & +2 \\
42 & 115 & not tabulated & no improved witness claimed & none \\
\end{longtable}

The literature review includes author-hosted sources and code, rather
than treating the newest table as complete. No source found in the
currently audited set supersedes these two comparisons. The completed
independent review supports the wording: ``Two exact constructions
improve the lower bounds in Ahmed--Snevily (2013); our literature audit
through 4 October 2026 found no prior equal or stronger construction.''
Its archive is \texttt{wave2/packing/few\_distance\_novelty/VERDICT.md}.
It is not a proof of exhaustive literature coverage. The withdrawn
September 2026 revision of arXiv:2510.09800 is excluded as a source of
established asymptotic theorems; its old mirror PDF remains accessible
{[}7{]}.

\subsection{3. Explicit constructions}\label{explicit-constructions}

For each column below, define S = \{(a,b) \ensuremath{\in} Z\ensuremath{{}^2} : Aa+Bb \ensuremath{\le} C for every
listed row\}. Thus the following ten inequalities alone define each
witness, including every interior lattice site.

\begin{longtable}[]{@{}rrrr@{}}
\toprule\noalign{}
A & B & C for S31 & C for S41 \\
\midrule\noalign{}
\endhead
\bottomrule\noalign{}
\endlastfoot
-2 & -1 & 9 & 10 \\
-1 & -2 & 9 & 11 \\
1 & -1 & 9 & 11 \\
1 & 0 & 5 & 7 \\
2 & 1 & 8 & 10 \\
1 & 1 & 5 & 6 \\
1 & 2 & 8 & 9 \\
0 & 1 & 5 & 6 \\
-1 & 1 & 8 & 9 \\
-1 & 0 & 5 & 6 \\
\end{longtable}

For S31, -5\ensuremath{\le}a\ensuremath{\le}5 and a-b\ensuremath{\le}9 imply b\ensuremath{\ge}-14, while b\ensuremath{\le}5. For S41, -6\ensuremath{\le}a\ensuremath{\le}7 and
a-b\ensuremath{\le}11 imply b\ensuremath{\ge}-17, while b\ensuremath{\le}6. Hence each polygon's complete integer
enumeration is contained in {[}-20,20{]}\ensuremath{{}^2}. There are no unexamined
points beyond the enumeration box.

Equivalently, for each listed a, choose every integer b between the
row's minimum and maximum, inclusive; choose no other rows. The
following intervals are obtained directly by substituting a into the ten
inequalities.

\subsubsection{S31}\label{s31}

\begin{longtable}[]{@{}rrrr@{}}
\toprule\noalign{}
a & minimum b & maximum b & sites \\
\midrule\noalign{}
\endhead
\bottomrule\noalign{}
\endlastfoot
-5 & 1 & 3 & 3 \\
-4 & -1 & 4 & 6 \\
-3 & -3 & 5 & 9 \\
-2 & -3 & 5 & 9 \\
-1 & -4 & 4 & 9 \\
0 & -4 & 4 & 9 \\
1 & -5 & 3 & 9 \\
2 & -5 & 3 & 9 \\
3 & -6 & 2 & 9 \\
4 & -5 & 0 & 6 \\
5 & -4 & -2 & 3 \\
\end{longtable}

The row lengths are 3,6,9,9,9,9,9,9,9,6,3, summing to 81. All sites
inside the ten-facet polygon are selected.

\subsubsection{S41}\label{s41}

\begin{longtable}[]{@{}rrrr@{}}
\toprule\noalign{}
a & minimum b & maximum b & sites \\
\midrule\noalign{}
\endhead
\bottomrule\noalign{}
\endlastfoot
-6 & 2 & 3 & 2 \\
-5 & 0 & 4 & 5 \\
-4 & -2 & 5 & 8 \\
-3 & -4 & 6 & 11 \\
-2 & -4 & 5 & 10 \\
-1 & -5 & 5 & 11 \\
0 & -5 & 4 & 10 \\
1 & -6 & 4 & 11 \\
2 & -6 & 3 & 10 \\
3 & -7 & 3 & 11 \\
4 & -7 & 2 & 10 \\
5 & -6 & 0 & 7 \\
6 & -5 & -2 & 4 \\
7 & -4 & -4 & 1 \\
\end{longtable}

The row lengths are 2,5,8,11,10,11,10,11,10,11,10,7,4,1, summing to 111.
All sites inside the corresponding polygon are selected.

\subsection{4. Complete elementary
certificate}\label{complete-elementary-certificate}

Let a row-interval set S have rows R\_a={[}L\_a,U\_a{]}\ensuremath{\cap}Z. For an
integer difference vector (u,v), its ordered-pair multiplicity is

N(u,v) = \ensuremath{\sum}\emph{a max(0, min(U\_a,U}\{a-u\}+v) - max(L\_a,L\_\{a-u\}+v)
+ 1),

where a term is zero if either row does not exist. This formula counts
the integers b for which both (a,b) and (a-u,b-v) lie in S: their
intervals intersect exactly as shown. For q\textgreater0 the number of
\textbf{unordered} pairs at squared distance q is

C(q) = (1/2) \ensuremath{\sum}\_\{u\ensuremath{{}^2}+uv+v\ensuremath{{}^2}=q\} N(u,v).

Only \textbar u\textbar\ensuremath{\le}a\_max-a\_min and
\textbar v\textbar\ensuremath{\le}b\_max-b\_min can contribute. These are finite
integer ranges. Negating (u,v) reverses each pair, explaining the factor
1/2. Zero differences are omitted. This proves the counting formula
without a solver or a floating-point operation.

For S31 use u\ensuremath{\in}{[}-10,10{]}, v\ensuremath{\in}{[}-11,11{]}. For S41 use
u,v\ensuremath{\in}{[}-13,13{]}. Substitution of the row intervals gives exactly the
positive multiplicities below; all other q have multiplicity zero.

\begin{longtable}[]{@{}rrr@{}}
\toprule\noalign{}
squared distance q & unordered pairs in S31 & unordered pairs in S41 \\
\midrule\noalign{}
\endhead
\bottomrule\noalign{}
\endlastfoot
1 & 208 & 291 \\
3 & 189 & 270 \\
4 & 175 & 252 \\
7 & 314 & 462 \\
9 & 144 & 216 \\
12 & 135 & 207 \\
13 & 254 & 390 \\
16 & 115 & 180 \\
19 & 212 & 342 \\
21 & 198 & 324 \\
25 & 88 & 147 \\
27 & 81 & 144 \\
28 & 158 & 276 \\
31 & 146 & 258 \\
36 & 63 & 117 \\
37 & 112 & 222 \\
39 & 108 & 216 \\
43 & 98 & 198 \\
48 & 36 & 81 \\
49 & 110 & 249 \\
52 & 64 & 156 \\
57 & 54 & 144 \\
61 & 40 & 114 \\
63 & 36 & 108 \\
64 & 19 & 60 \\
67 & 28 & 102 \\
73 & 16 & 90 \\
75 & 9 & 36 \\
76 & 16 & 72 \\
79 & 10 & 66 \\
81 & 0 & 36 \\
84 & 0 & 54 \\
91 & 4 & 84 \\
93 & 0 & 36 \\
97 & 0 & 30 \\
100 & 0 & 12 \\
103 & 0 & 18 \\
108 & 0 & 9 \\
109 & 0 & 18 \\
112 & 0 & 12 \\
127 & 0 & 6 \\
\end{longtable}

There are 31 positive entries in the S31 column and 41 in the S41
column. Their sums are respectively 3240=81·80/2 and 6105=111·110/2. The
finite difference ranges cover every possible pair, so the table is
complete. Every listed positive value is realized; no other value
occurs. The injectivity of \ensuremath{\varphi} and the positivity identity for Q now prove
both theorems.

A second computation enumerates unordered point pairs and uses the
Cartesian numerator (2(a-c)+b-d)\ensuremath{{}^2}+3(b-d)\ensuremath{{}^2}, dividing by 4. Agreement with
the row-correlation calculation is archived in
\texttt{wave2/paper/evidence.json}. This cross-check is manuscript
reproducibility; the independent certification reported below uses
separately implemented checkers.

The full squared-distance set for S31 is

1, 3, 4, 7, 9, 12, 13, 16, 19, 21, 25, 27, 28, 31, 36, 37, 39, 43, 48,
49, 52, 57, 61, 63, 64, 67, 73, 75, 76, 79, 91.

For S41 it is

1, 3, 4, 7, 9, 12, 13, 16, 19, 21, 25, 27, 28, 31, 36, 37, 39, 43, 48,
49, 52, 57, 61, 63, 64, 67, 73, 75, 76, 79, 81, 84, 91, 93, 97, 100,
103, 108, 109, 112, 127.

\subsection{5. Geometry and limitations of the structural
interpretation}\label{geometry-and-limitations-of-the-structural-interpretation}

S31 is a filled convex decagon. Among lattice rotations and reflections
followed by translation, its only symmetries are identity and
(a,b)\ensuremath{\mapsto}(-a,a+b), a Euclidean reflection. This follows by transforming the
complete finite set under the twelve triangular-lattice dihedral maps
and checking the only possible aligning translations. A single
asymmetric witness does not contradict a conjecture about the symmetry
of a globally optimal configuration: no global optimality is proved.

S41 has centroid (1/3,-2/3) in lattice coordinates. The map
(a,b)\ensuremath{\mapsto}(-a-b,a-1) is a 120-degree rotation about this centroid, and
(a,b)\ensuremath{\mapsto}(b+1,a-1) is a reflection fixing it. The corresponding six D3 maps
preserve the entire set. Accordingly all nonzero-distance multiplicities
in S41 are multiples of three. This observation is explanatory, not part
of the existence proof.

S31 uses squared norm 91 while avoiding 81 and 84; it therefore is not
restricted to the 31 shortest triangular-lattice norms. S41 similarly
avoids 111,117,121 while using 112 and 127. These are concrete alphabet
choices. Their value is the certified cardinality improvement at the two
specified indices; the principle of trading distance classes predates
this work.

The known symmetric saw-blade families B\_s,r,t in {[}2{]} do not
represent these two witnesses up to similarity: cardinality restricts
the standard families to parameter cases whose core distances include
norms absent from our witnesses. The finite arithmetic exclusion is
archived separately. A necessary search for a polygonal P\_a,r1,r2 core
with individually attached unit-triangle sites also found no match
within the tested parameter ranges. That computation is not a proof
excluding every interpretation of arbitrary blades, and neither
exclusion by itself establishes novelty against all prior work.

\subsubsection{An exact parameter construction for
S41}\label{an-exact-parameter-construction-for-s41}

For integers s\ensuremath{\ge}1 and t\ensuremath{\ge}0, let P(s,t) be the polygon with successive
vertices

\begin{verbatim}
(0,0), (s,-2s), (3s,-3s), (3s+t,-3s),
(4s+t,-2s), (3s+t,0), (3s,t), (s,s+t), (0,t).
\end{verbatim}

Define F(s,t)=P(s,t)\ensuremath{\cap}Z\ensuremath{{}^2}, selecting every lattice site. For
t\textgreater0 the edge vectors are

\begin{verbatim}
(s,-2s), (2s,-s), (t,0), (s,s), (-s,2s),
(-t,t), (-2s,s), (-s,-s), (0,-t).
\end{verbatim}

They sum to zero and their Euclidean directions turn by 60°,30°,30°,
repeated three times. The closed polygon is strictly convex. At t=0 the
three zero edges collapse to a regular hexagon. The shoelace sum and the
sum of edge-coordinate gcds are respectively

2A = 18s\ensuremath{{}^2}+12st+t\ensuremath{{}^2}, B = 6s+3t.

Here A is area in lattice-coordinate space, where Z\ensuremath{{}^2} has unit
fundamental cell; it is not area in the physical \ensuremath{\varphi} embedding. Pick's
theorem gives the following exact cardinality identity for every allowed
s,t:

\textbf{Proposition 3.} \textbar F(s,t)\textbar{} =
9s\ensuremath{{}^2}+6st+t(t+3)/2+3s+1.

\textbf{Proof.} The listed closed vertices give the displayed shoelace
sum. Six edges have coordinate gcd s, and three have gcd t, giving B.
Counting interior and boundary sites, Pick's theorem gives
\textbar F\textbar=A+B/2+1. Substitution yields the identity. At t=0,
the same computation applies after collapsing zero-length edges.
\(\square\)

The construction is D3-symmetric for t\textgreater0 about
c=(2s+t/3,-s+t/3). The centered 120-degree rotation is
(a,b)\ensuremath{\mapsto}(-a-b,a)+(3s+t,-3s); the centered reflection is
(a,b)\ensuremath{\mapsto}(b,a)+(3s,-3s). They preserve the listed vertices and Z\ensuremath{{}^2}, so they
preserve every filled lattice site. The three distinguished t-edges have
Euclidean length t and the other six length s\ensuremath{\surd}3. These lengths are
unequal for positive integers s,t, restricting any extra symmetry to
that of the triangle formed by the distinguished-edge midpoints; the
full group is D3.

At (s,t)=(3,1), the identity gives 111 sites: doubled coordinate area is
199 and B=21. Direct translation shows F(3,1)=S41+(6,-2). The pair
certificate in Section 4 proves its 41 distances. The general
cardinality identity is an algebraic consequence of the explicit
polygon, not an extrapolation from a finite batch. It is not a claimed
formula for the number of distances. The symbolic area and D3
identities, and the translated 111-point polygon equality, are
separately checked by
\texttt{wave2/paper/verify\_family\_proposition.py}; its result is
archived in \texttt{family\_proposition\_certificate.json}. The executed
72-case parameter batch s\ensuremath{\in}\{1,\ldots,8\},t\ensuremath{\in}\{0,\ldots,8\} checked the
counts; only the 111/41 case improved the audited comparison table in
that batch. No new asymptotic mechanism or generally improving infinite
family is claimed.

\subsubsection{5.2. Why squared distance 111
disappears}\label{why-squared-distance-111-disappears}

For F(s,1), the listed vertices give the exact widths 0\ensuremath{\le}a\ensuremath{\le}4s+1 and
0\ensuremath{\le}2a+b\ensuremath{\le}6s+2. Every point difference therefore satisfies
\textbar a\textbar\ensuremath{\le}4s+1 and \textbar2a+b\textbar\ensuremath{\le}6s+2. The D3 symmetry
of the polygon and central symmetry of its difference body generate D6.
Applying these symmetries gives

\textbar a\textbar, \textbar b\textbar, \textbar a+b\textbar{} \ensuremath{\le} 4s+1,

\textbar2a+b\textbar, \textbar a+2b\textbar, \textbar a-b\textbar{} \ensuremath{\le}
6s+2.

The vector (3s+1,1) violates the latter bound by one and has squared
norm 9s\ensuremath{{}^2}+9s+3=N(s,1). Its whole twelve-element dihedral orbit is
excluded. Other representations of this same norm can survive: this is
not a universal absence theorem for that norm.

At s=3, norm111 has precisely that excluded orbit. Completing a square
bounds every representation by
\textbar a\textbar,\textbar b\textbar\textless13; exact enumeration in
that finite box proves the orbit is complete. More generally, enumerate
the nonzero integer vectors satisfying the six displayed caps at s=3.
There are420vectors and41norms, maximum127. Since every actual point
difference belongs to this envelope, it independently proves at
most41distances. The full pair histogram proves all41occur. The palette
is precisely the first41positive triangular-lattice norms
with111replaced by127; among the45representable norms up to127, it
omits111,117,121,124. The support-width argument explains this
particular improvement; the broader strategy of substituting distance
classes is prior art. Independent analytic replay is
\texttt{wave2/final\_review/orbit\_review.py}.

The shared half-plane representation suggests useful construction
programs. Scaling or changing its supports produces many exact point
sets, but the present evidence establishes neither a uniformly improving
infinite family nor an asymptotic theorem. Cases with no audited
comparator are not described as improvements.

\subsection{6. Discovery, independent checking, and
reproducibility}\label{discovery-independent-checking-and-reproducibility}

The 81-point witness arose from joint Boolean selection of lattice sites
and distance classes in a radius-6 hexagonal window. For every pair,
selecting both sites forces its squared-norm variable; at most 31 norm
variables may be selected. Valid matching cuts accelerate the search.
Seed 42102 returned 81 points in 4.143 seconds and nine search nodes.
That runtime includes the logged finite instance, not the entire earlier
research portfolio. A radius-7 continuation with a fixed 37-site
interior reproduced 81, without improving it. Solver optimality in a
finite or partly fixed model is never used as planar optimality.

The 81-point support normals provided the starting point for a separate
enumeration of support offsets. The delta-2 run evaluated 9,765,625
support tuples in 54.379 seconds and produced the 111-point witness.
This is independent construction development from a shared lead, not an
independent rediscovery from an untouched problem statement. A
subsequent canonical search evaluated 5,764,801 tuples in 31.9019
seconds. These are recorded batch timings, not the exact time to the
first 111-point observation.

An independent integer checker and the separate
Cartesian/common-denominator verifier agreed on the original 81-point
witness. For the 111-point witness, the same independent Cartesian
verifier recomputed all 6105 pairs without importing the
support-enumeration search code. Candidate geometry, historical source
values, and novelty are reviewed separately. No independent checker
accepts discovery metadata as a mathematical bound.

Reproduce both frozen witnesses, their half-plane/row agreement, every
pair multiplicity, and artifact hashes with standard Python:

\begin{verbatim}
python3 research/geometric-extremal/results/reproduce.py
python3 research/geometric-extremal/wave2/paper/verify_witnesses.py
\end{verbatim}

No optimization package, internet access, numerical tolerance, or random
seed is needed for this command. The verifier regenerates coordinates
from the manuscript row definitions and compares them with the frozen
files. The complete source-and-search archive also retains model texts,
seeds, prospective registrations, interrupted experiments, negative
results, and explicit baseline controls.

Witness hashes are listed below; file paths are relative to the research
directory.

GE-LOWER-G31:
\texttt{candidates/wave2\_joint\_palette/k31\_n81\_r6\_seed42102.json}

\begin{verbatim}
a85a2a990531ebd5b40165eacb0f23c84bac6676404cb858f7b5c2c63b83d94c
\end{verbatim}

GE-LOWER-G41:
\texttt{candidates/wave2\_family/support\_k41\_n111\_code7324998.json}

\begin{verbatim}
2d1216002f65a2c2cccbaa30299ea4f26c21c6f493fe3ad0e164de11ad0aa3b1
\end{verbatim}

The mathematical definitions and integer proof do not depend on those
hashes; the hashes identify the exact files used in the reported checks.
A corrupted or changed file is rejected by the frozen reproduction
command.

\subsection{References}\label{references}

{[}1{]} P. Erdős and P. Fishburn. \emph{Maximum planar sets determining
k distances.} Discrete Mathematics 160 (1996), 115--125. Bibliographic
details cross-checked in {[}4{]}; original paper remains the
foundational problem reference.

{[}2{]} T. Ahmed and H. Snevily.
\href{https://www.combinatorics.org/ojs/index.php/eljc/article/download/v20i4p33/pdf/}{\emph{Sparse
Distance Sets in the Triangular Lattice}}. Electronic Journal of
Combinatorics 20(4) (2013), P33. Printed page 5 gives the target lower
bounds. Author publication index:
https://tanbir.github.io/publications.html.

{[}3{]} T. Ahmed and D. Wildstrom.
\href{https://tanbir.github.io/files/preprints/ahmed_wildstrom_distance_sets.pdf}{\emph{On
Distance Sets in the Triangular Lattice}}. Bulletin of the Institute of
Combinatorics and its Applications 75 (2015), 118--127. Proposition 2.1
supplies the boundary-difference result.

{[}4{]} (a) V. Balaji, O. Edwards, A. M. Loftin, S. Mcharo, A. Rice, and
B. Tsegaye.
\href{https://math.colgate.edu/~integers/u86/u86.pdf}{\emph{Lattice
Configurations Determining Few Distances}}. Integers 20 (2020), A86;
published 19 October 2020. This journal PDF has six authors.

\begin{enumerate}
\def\labelenumi{(\alph{enumi})}
\setcounter{enumi}{1}
\tightlist
\item
  V. Balaji, O. Edwards, A. M. Loftin, S. Mcharo, L. Phillips, A. Rice,
  and B. Tsegaye. \href{https://arxiv.org/abs/1911.11688}{\emph{Lattice
  Configurations Determining Few Distances}}. arXiv:1911.11688v3, 7 July
  2023; initial preprint 26 November 2019. This revised preprint has
  seven authors.
  \href{https://alexricemath.com/wp-content/uploads/2023/07/LatticeDraftRevisedNew.pdf}{Author-hosted
  current PDF}.
\end{enumerate}

{[}5{]} F. Szöllősi and P. R. J. Östergård.
\href{https://www.combinatorics.org/ojs/index.php/eljc/article/view/v27i1p23}{\emph{Constructions
of Maximum Few-Distance Sets in Euclidean Spaces}}. Electronic Journal
of Combinatorics 27(1) (2020), P1.23. DOI 10.37236/8565; published 24
January 2020.

{[}6{]} L.-R. Bao and W.-H. Yu.
\href{https://arxiv.org/html/2509.00880}{\emph{Constructions of Large
m-Distance Sets on Triangular Lattice}}. arXiv:2509.00880v1, submitted
31 August 2025.
\href{https://github.com/BaoCoder613/triangular-m-distance}{Author
code}.

{[}7{]} L. Wang. \href{https://arxiv.org/abs/2510.09800}{\emph{On
Few-Distance Sets in the Plane}}. arXiv:2510.09800, v2 withdrawn 9
September 2026, with the stated reason of an error in the asymptotic
proof. Cited solely to document source status, not to support a theorem.

Primary web sources were checked on 4 October 2026. The completed
independent novelty review is archived at
\texttt{wave2/packing/few\_distance\_novelty/VERDICT.md}; its frozen
SHA256 is:

\begin{verbatim}
93c22b77200f6e0209fea172b6e065d7edd0f89404184947a1ac5a17c04f328a
\end{verbatim}

Its scoped conclusion supports improvements over audited published
bounds, while rejecting global-optimality and unconditional
worldwide-record claims.

\end{document}
